% ==========================================================================
%  IA na Sociedade — cabeçalho LaTeX (pdflatex)
%  Programa CiberExt 26-29 · FEELT38103
%  Adaptado de "AI in Society" (Universidade de Helsinque / Una Europa)
%  Licença: CC BY-NC-SA 4.0
% --------------------------------------------------------------------------
%  Espelha a identidade do ia-style.css: papel quente, tinta ardósia e a
%  faixa tricolor dos três eixos do curso (conceitos / aplicações / riscos).
% ==========================================================================

% ---------- Fontes ----------
% A serifada vem do build.sh, que passa  -V fontfamily=charter  ao pandoc.
% Aqui só acrescentamos a sem-serifa; ambas com fallback, porque nem toda
% instalação de LaTeX traz os mesmos pacotes.
\IfFileExists{helvet.sty}{\usepackage[scaled=0.90]{helvet}}{}
\usepackage{textcomp}
\IfFileExists{microtype.sty}{\usepackage{microtype}}{}

% ---------- Pacotes de layout ----------
\usepackage{xcolor}
\usepackage{tcolorbox}
\tcbuselibrary{skins,breakable}
\usepackage{titlesec}
\usepackage{fancyhdr}
\usepackage{enumitem}
\usepackage{amssymb}
\usepackage{booktabs}
\usepackage{array}
\usepackage{colortbl}
\usepackage{ragged2e}
\usepackage{fvextra}
\usepackage{newunicodechar}

% ---------- Caracteres Unicode que o pdflatex não conhece ----------
\newunicodechar{✓}{{\color{iaAplic}\checkmark}}
\newunicodechar{✅}{{\color{iaAplic}\checkmark}}
\newunicodechar{→}{\ensuremath{\rightarrow}}
\newunicodechar{←}{\ensuremath{\leftarrow}}
\newunicodechar{↔}{\ensuremath{\leftrightarrow}}
\newunicodechar{–}{--}
\newunicodechar{—}{---}
\newunicodechar{€}{\texteuro}
\newunicodechar{£}{\textsterling}
\newunicodechar{×}{\ensuremath{\times}}
\newunicodechar{≈}{\ensuremath{\approx}}
\newunicodechar{≥}{\ensuremath{\geq}}
\newunicodechar{≤}{\ensuremath{\leq}}
\newunicodechar{•}{\textbullet}
\newunicodechar{″}{''}
\newunicodechar{′}{'}

% ==========================================================================
%  PALETA — os mesmos valores do ia-style.css
% ==========================================================================
\definecolor{iaConceitos}{HTML}{24505C}   % petróleo — títulos, eixo 1
\definecolor{iaAplic}{HTML}{16746C}       % verdete  — subtítulos, eixo 2
\definecolor{iaRiscos}{HTML}{B45F14}      % açafrão  — destaques, eixo 3
\definecolor{iaAmeixa}{HTML}{77345F}      % links internos, termos técnicos
\definecolor{iaTinta}{HTML}{1F2933}       % corpo do texto
\definecolor{iaTintaFraca}{HTML}{5C6670}
\definecolor{iaRegra}{HTML}{E4DBCD}
\definecolor{iaBloco}{HTML}{F3EEE4}
\definecolor{iaNota}{HTML}{EAF1F0}
\definecolor{iaCaso}{HTML}{FCF3E6}
\definecolor{iaFilos}{HTML}{F7EFF4}

\color{iaTinta}

% ==========================================================================
%  ASSINATURA — faixa tricolor dos três eixos
% ==========================================================================
\newcommand{\faixaEixos}[1][\linewidth]{%
  \noindent
  \textcolor{iaConceitos}{\rule{0.3333#1}{4pt}}%
  \textcolor{iaAplic}{\rule{0.3333#1}{4pt}}%
  \textcolor{iaRiscos}{\rule{0.3333#1}{4pt}}%
  \par}

% ---------- Página de título ----------
\makeatletter
\newcommand{\subtitle}[1]{\gdef\@subtitle{#1}}
\providecommand{\@subtitle}{}
\def\@maketitle{%
  \newpage\null\vskip 1.5em%
  \faixaEixos
  \vskip 2.2em
  {\raggedright
    {\sffamily\bfseries\huge\color{iaConceitos}\@title\par}%
    \vskip 0.7em%
    {\sffamily\large\color{iaTintaFraca}\@subtitle\par}%
    \vskip 2em%
    {\color{iaRegra}\rule{\linewidth}{1pt}}%
    \vskip 1em%
    {\sffamily\small\color{iaTinta}\@author\par}%
    \vskip 0.5em%
    {\sffamily\footnotesize\color{iaTintaFraca}\@date\par}%
  }\par\vskip 2.5em}
\makeatother

% ==========================================================================
%  TÍTULOS DE SEÇÃO
% ==========================================================================
\titleformat{\section}[block]
  {\sffamily\Large\bfseries\color{iaConceitos}}
  {\thesection}{0.7em}{}
  [{\vspace{2pt}\color{iaRegra}\titlerule[1.2pt]}]
\titleformat{\subsection}
  {\sffamily\large\bfseries\color{iaAplic}}
  {\thesubsection}{0.6em}{}
\titleformat{\subsubsection}
  {\sffamily\normalsize\bfseries\color{iaRiscos}}
  {\thesubsubsection}{0.5em}{}
\titleformat{\paragraph}[runin]
  {\sffamily\small\bfseries\color{iaRiscos}}
  {}{0em}{}

\titlespacing*{\section}{0pt}{2.6ex plus 1ex minus .2ex}{1.3ex plus .2ex}
\titlespacing*{\subsection}{0pt}{2.2ex plus 1ex minus .2ex}{1ex plus .2ex}

% ==========================================================================
%  CABEÇALHO E RODAPÉ
% ==========================================================================
\setlength{\headheight}{24pt}
\renewcommand{\sectionmark}[1]{\markboth{#1}{}}
\pagestyle{fancy}
\fancyhf{}
\renewcommand{\headrulewidth}{0.6pt}
\renewcommand{\footrulewidth}{0pt}
\fancyhead[L]{\sffamily\footnotesize\color{iaAplic}Ética da IA}
\fancyhead[R]{\sffamily\footnotesize\color{iaAplic}\nouppercase{\leftmark}}
\fancyfoot[C]{\sffamily\footnotesize\color{iaConceitos}\thepage}
\renewcommand{\headrule}{%
  \hbox to\headwidth{\color{iaRegra}\leaders\hrule height \headrulewidth\hfill}}

% ==========================================================================
%  TEXTO CORRIDO
% ==========================================================================
\linespread{1.08}
\setlength{\parskip}{0.55em}
\setlength{\parindent}{0pt}
\emergencystretch=3em
\hbadness=10000

% ---------- Termos técnicos / código inline ----------
\let\oldtexttt\texttt
\renewcommand{\texttt}[1]{{\color{iaAmeixa}\small\ttfamily #1}}

% ---------- Blocos verbatim ----------
\fvset{breaklines=true,breakanywhere=true,fontsize=\small}

% ---------- Blocos de código (ambiente Shaded do pandoc) ----------
% O pandoc só pré-define o ambiente Shaded quando o markdown tem bloco de
% código; nos capítulos sem código ele não existe, e \renewenvironment
% incondicional falharia com "Environment Shaded undefined". \@ifundefined
% escolhe \newenvironment ou \renewenvironment conforme o caso.
\makeatletter
\@ifundefined{Shaded}{\newenvironment{Shaded}}{\renewenvironment{Shaded}}
  {\begin{tcolorbox}[
      breakable, enhanced jigsaw, rounded corners,
      arc=3pt, boxrule=0pt, frame hidden,
      colback=iaBloco,
      borderline west={3pt}{0pt}{iaAmeixa},
      left=10pt, right=8pt, top=6pt, bottom=6pt,
      before skip=9pt, after skip=9pt]}
  {\end{tcolorbox}}
\makeatother

% ---------- Citações: caixa de caso ----------
\renewenvironment{quote}
  {\begin{tcolorbox}[
      breakable, enhanced jigsaw, rounded corners,
      arc=3pt, boxrule=0pt, frame hidden,
      colback=iaCaso,
      borderline west={3pt}{0pt}{iaRiscos},
      left=12pt, right=10pt, top=9pt, bottom=9pt,
      before skip=9pt, after skip=9pt]}
  {\end{tcolorbox}}

% ==========================================================================
%  BLOCOS DE APOIO
%  No markdown, escreva:   ::: nota   ...texto...   :::
%  O pandoc converte fenced_divs para os ambientes abaixo.
% ==========================================================================
\newtcolorbox{iabox}[3]{
  breakable, enhanced jigsaw, rounded corners,
  arc=3pt, boxrule=0pt, frame hidden,
  colback=#2,
  borderline west={3pt}{0pt}{#1},
  left=12pt, right=10pt, top=9pt, bottom=9pt,
  before skip=11pt, after skip=11pt,
  before upper={\setlength{\parskip}{0.5em}\setlength{\parindent}{0pt}},
  title={\sffamily\bfseries\footnotesize\textcolor{#1}{\MakeUppercase{#3}}},
  attach title to upper={\par\vspace{5pt}},
  fonttitle=\normalfont}

% Cada ambiente aceita um titulo opcional entre colchetes, vindo do
% atributo data-titulo do markdown:   \begin{filosofia}[Titulo] ... \end{filosofia}
% Os quatro primeiros espelham os icones do material original.
\newenvironment{filosofia}[1][Conceito filosofico]{\begin{iabox}{iaAmeixa}{iaFilos}{#1}}{\end{iabox}}
\newenvironment{tecnica}[1][Nota tecnica]{\begin{iabox}{iaConceitos}{iaNota}{#1}}{\end{iabox}}
\newenvironment{contexto}[1][Contexto]{\begin{iabox}{iaAplic}{iaBloco}{#1}}{\end{iabox}}
\newenvironment{definicao}[1][Definicao]{\begin{iabox}{iaRiscos}{iaCaso}{#1}}{\end{iabox}}

\newenvironment{nota}[1][Nota]{\begin{iabox}{iaConceitos}{iaNota}{#1}}{\end{iabox}}
\newenvironment{caso}[1][Caso]{\begin{iabox}{iaRiscos}{iaCaso}{#1}}{\end{iabox}}
\newenvironment{reflexao}[1][Para pensar]{\begin{iabox}{iaAplic}{iaBloco}{#1}}{\end{iabox}}
\newenvironment{licenca}
  {\par\vspace{2em}\faixaEixos\par\vspace{0.8em}
   \begingroup\sffamily\scriptsize\color{iaTintaFraca}}
  {\endgroup\par}

% ==========================================================================
%  LISTAS
% ==========================================================================
\setlist[itemize]{leftmargin=1.4em, itemsep=2pt, topsep=4pt, parsep=0pt}
\setlist[enumerate]{leftmargin=1.7em, itemsep=2pt, topsep=4pt, parsep=0pt}
\renewcommand{\labelitemi}{\textcolor{iaRiscos}{\textbullet}}
\renewcommand{\labelitemii}{\textcolor{iaAplic}{--}}

% ==========================================================================
%  TABELAS
% ==========================================================================
\renewcommand{\arraystretch}{1.3}
\arrayrulecolor{iaRegra}

% ==========================================================================
%  RÓTULOS EM PORTUGUÊS
%  Definidos aqui (e não via -V na linha de comando) porque acentos passados
%  como argumento dependem do locale do terminal e podem se corromper.
% ==========================================================================
\renewcommand{\contentsname}{Sumário}
\renewcommand{\refname}{Referências}
\renewcommand{\figurename}{Figura}
\renewcommand{\tablename}{Tabela}
